\chapter{Kinematics: describing motion}\label{ch:kinematics}
\index{kinematics}

Mechanics begins with description, not explanation: before we ask \emph{why}
things move (chapter~2), we need language precise enough to say
\emph{how} they move. That language is calculus applied to a position function
$x(t)$, and the payoff for learning it carefully is that the hardest parts of
physics --- momentum, orbits, wave motion --- reuse exactly these ideas. This
chapter derives the constant-acceleration formulas from the definitions (never
quotes them), treats projectile motion as two independent one-dimensional
problems, and builds the centripetal-acceleration formula from a picture of
velocity differences. Engineering practice shows up at the end: inertial
navigation is this chapter run continuously.

\section{Position, displacement, and velocity}\label{sec:kin-pos}

Fix a line and an origin. A point particle's \emph{position}\index{position}
$x(t)$ is a signed distance from the origin; the sign encodes direction. The
\emph{displacement}\index{displacement} over $[t_1,t_2]$ is
$\Delta x=x(t_2)-x(t_1)$ --- which is emphatically not the distance travelled:
run around a track once and the displacement is zero while the distance is a
lap. The \emph{average velocity}\index{average velocity} is the slope of the
secant,
\[
  \bar v=\frac{\Delta x}{\Delta t},
\]
and the \emph{instantaneous velocity} is what the average converges to as the
interval shrinks:
\[
  v(t)=\lim_{\Delta t\to0}\frac{x(t+\Delta t)-x(t)}{\Delta t}
      =\frac{dx}{dt}.
\]
Velocity is a slope on the $x(t)$ graph. Speed is its magnitude --- ``60 mph
east'' is a velocity, ``60 mph'' is a speed. Sign matters: in one dimension
$v>0$ means moving toward $+x$.

\begin{rulebox}[Reading an $x(t)$ graph]
Slope of the curve at $t$ = instantaneous velocity. Slope of the secant over
$[t_1,t_2]$ = average velocity. A horizontal tangent = momentarily at rest ---
and that is all it means; what happens next depends on the acceleration.
\end{rulebox}

\begin{exercise}[easy]
A particle sits at $x(t)=3t^{2}-12t$. Find its average velocity over
$[0,4]$ and its instantaneous velocity at $t=2$. (AP-C)
\end{exercise}

\section{Acceleration and the calculus loop}\label{sec:kin-acc}

Acceleration is velocity's velocity:
\[
  a(t)=\frac{dv}{dt}=\frac{d^{2}x}{dt^{2}},
\]
and figure~\ref{fig:triple} shows all three functions for a thrown object.
Read them vertically: when $v$ crosses zero the object is at the top of its
climb, and $a$ --- a constant $-9.8\,\mathrm{m/s^{2}}$ --- does not care.

\begin{figure}[htbp]\centering
% x(t), v(t), a(t) aligned for x = 20t - 4.9 t^2 (thrown object)
\begin{tikzpicture}
\begin{axis}[name=px, width=0.86\linewidth, height=3.4cm,
  xlabel={}, ylabel={$x$ (m)}, ymin=0, ymax=22, xmin=0, xmax=4.2,
  grid=major, grid style={gray!20}, xtick=\empty]
\addplot[domain=0:4.1,samples=60,thick,accent]{20*x-4.9*x^2};
\end{axis}
\begin{axis}[name=pv, at={(px.below south west)}, anchor=above north west,
  yshift=-2mm, width=0.86\linewidth, height=3.4cm,
  xlabel={}, ylabel={$v$ (m/s)}, ymin=-22, ymax=22, xmin=0, xmax=4.2,
  grid=major, grid style={gray!20}, xtick=\empty]
\addplot[domain=0:4.1,samples=60,thick,accent]{20-9.8*x};
\draw[dashed,gray] (axis cs:2.04,-22) -- (axis cs:2.04,22);
\end{axis}
\begin{axis}[name=pa, at={(pv.below south west)}, anchor=above north west,
  yshift=-2mm, width=0.86\linewidth, height=2.6cm,
  xlabel={$t$ (s)}, ylabel={$a$ (m/s$^2$)}, ymin=-15, ymax=8, xmin=0, xmax=4.2,
  grid=major, grid style={gray!20}]
\addplot[domain=0:4.1,samples=2,thick,accent]{-9.8};
\draw[dashed,gray] (axis cs:2.04,-15) -- (axis cs:2.04,8);
\end{axis}
\node[font=\footnotesize,rotate=90] at ($(px.west)+(-0.55,0)$) {peak};
\end{tikzpicture}

\caption{$x(t)=20t-4.9t^{2}$ with its first two derivatives, drawn on one
time axis. Notice the three-way alignment: the peak of $x$ sits exactly above
the zero of $v$, while $a$ stays fixed at $-9.8\,\mathrm{m/s^{2}}$ the whole
flight. A horizontal tangent in $x$ means ``stopped \emph{then},'' never
``stopped forever.''}
\label{fig:triple}
\end{figure}

The loop of figure~\ref{fig:loop} is the chapter's engine: differentiate to go
from $x$ to $v$ to $a$; integrate to come back:
\[
  v(t)=v(t_0)+\int_{t_0}^{t}a(t')\,dt',\qquad
  x(t)=x(t_0)+\int_{t_0}^{t}v(t')\,dt'.
\]
Kinematics is nothing more --- and nothing less --- than choosing a direction
around that loop.

\begin{figure}[htbp]\centering
% the calculus loop: x -> v -> a by d/dt, back by integrals
\begin{tikzpicture}[scale=1.0,>=stealth,
  box/.style={draw,thick,minimum width=2.2cm,minimum height=0.9cm,font=\small}]
  \node[box] (x) at (0,0) {position $x(t)$};
  \node[box] (v) at (4.0,0) {velocity $v(t)$};
  \node[box] (a) at (8.0,0) {acceleration $a(t)$};
  \draw[->,very thick] (x) -- (v) node[midway,above,font=\small] {$d/dt$};
  \draw[->,very thick] (v) -- (a) node[midway,above,font=\small] {$d/dt$};
  \draw[->,very thick,dashed,accent] (a.south) .. controls (5.5,-1.6) .. (v.south)
    node[midway,below,font=\small] {$v(t)=v_0+\int a\,dt$};
  \draw[->,very thick,dashed,accent] (v.south) .. controls (1.5,-1.6) .. (x.south)
    node[midway,below,font=\small] {$x(t)=x_0+\int v\,dt$};
  \node[font=\footnotesize,align=center] at (4.0,1.1) {differentiate going right;\\integrate coming back};
\end{tikzpicture}

\caption{The calculus loop. Notice that going right loses information (the
constants $x_0,v_0$ are invisible in $a(t)$), so coming back needs initial
conditions: two numbers must be supplied before any $x(t)$ is unique.}
\label{fig:loop}
\end{figure}

\begin{example}[Slopes off a measured $x(t)$]\label{ex:dataderiv}
A radar logs a dragster's position every half second:
\begin{center}
\begin{tabular}{l|ccccccc}
$t$ (s) & 0 & 0.5 & 1.0 & 1.5 & 2.0 & 2.5 & 3.0\\\hline
$x$ (m) & 0 & 3.1 & 9.8 & 22.0 & 39.5 & 62.4 & 90.7
\end{tabular}
\end{center}
Estimate $v(1.0)$ and $a(1.0)$ without a formula for $x(t)$.
\par\smallskip
\textbf{Strategy.} With data, derivatives are differences: centre the secant
for $v$ (symmetric difference quotients are second-order accurate), then
difference the velocities for $a$.
\begin{solution}
$v(1.0)\approx\frac{x(1.5)-x(0.5)}{1.0}=\frac{22.0-3.1}{1.0}=
18.9\,\mathrm{m/s}$. For $a$: estimate $v(0.5)\approx9.8/1.0=9.8$ (from
$x(1.0)-x(0)$ over 1\,s $\Rightarrow$ centre it at 0.5) and
$v(1.5)\approx(39.5-9.8)/1.0=29.7$; then
$a(1.0)\approx(29.7-9.8)/1.0=19.9\,\mathrm{m/s^{2}}\approx2g$.
\checkmark\ Same calculus, no closed form needed --- this is exactly how IMU
and encoder data get turned into velocities.
\end{solution}
\end{example}

\section{Constant acceleration, derived}\label{sec:kin-consta}

When $a$ is constant the loop collapses to three formulas --- derived here so
their assumptions are visible. Integrating $a$ once,
$v(t)=v_0+at$; integrating again,
\begin{equation}
  x(t)=x_0+v_0 t+\tfrac12 at^{2}.\label{eq:consta1}
\end{equation}
Eliminating $t$ between them (solve $v=v_0+at$ for $t$, substitute into
\eqref{eq:consta1}) gives the time-free relation
\begin{equation}
  v^{2}=v_0^{2}+2a\,(x-x_0).\label{eq:consta2}
\end{equation}
That is the whole subject of ``constant-acceleration kinematics.'' Everything
else --- free fall, braking distance, the meeting-time problem --- is one of
these two equations with different letters.

\begin{example}[Braking distance and the $v^{2}$ trap]\label{ex:brake}
A car at $100\,\mathrm{km/h}$ brakes at a steady
$a=-8\,\mathrm{m/s^{2}}$ (good tires, dry road). How far does it travel before
stopping --- and how far if the driver needs $0.5\,\mathrm{s}$ to react first?
\par\smallskip
\textbf{Strategy.} No time is asked or given: equation \eqref{eq:consta2}
answers it directly. Convert units before plugging; add the reaction-time
distance afterwards (during it, $a=0$).
\begin{solution}
$100\,\mathrm{km/h}=27.8\,\mathrm{m/s}$. From \eqref{eq:consta2},
$0=v_0^{2}+2a\,\Delta x$, so
$\Delta x=\frac{v_0^{2}}{2|a|}=\frac{27.8^{2}}{16}=48.3\,\mathrm{m}$.
Reaction adds $v_0\Delta t=27.8\times0.5=13.9\,\mathrm{m}$ --- total
$62.2\,\mathrm{m}$, a third of it before the brakes touch. \checkmark\
Notice the $v^{2}$: at $140\,\mathrm{km/h}$ the braking part alone grows to
$94\,\mathrm{m}$, not $67$ --- doubling speed quadruples the stop.
\end{solution}
\end{example}

\begin{exercise}[medium]
Two trains $400\,\mathrm{m}$ apart approach at $20\,\mathrm{m/s}$ each. When
does a bee flying between them at $30\,\mathrm{m/s}$ get squashed --- and how
far does the bee travel? Solve without summing the series. (F=ma)
\end{exercise}

\section{Projectile motion: two problems at once}\label{sec:kin-proj}

With gravity the only acceleration, the vector equation
$\mathbf a=(0,-g)$ splits: horizontal motion is constant-velocity,
vertical is constant-acceleration, and they share only the clock:
\[
  x(t)=v_{0x}t,\qquad
  y(t)=v_{0y}t-\tfrac12 gt^{2}.
\]
Independence is the deep point --- a dropped ball and a fired one land
together --- and figure~\ref{fig:proj} shows the trajectory with its velocity
components at three instants.

\begin{figure}[htbp]\centering
% projectile with velocity components at launch, apex, descent
\begin{tikzpicture}
\begin{axis}[width=0.9\linewidth,height=6.0cm,
  xlabel={$x$ (m)}, ylabel={$y$ (m)}, xmin=0, xmax=44, ymin=0, ymax=13,
  grid=major, grid style={gray!20}, axis equal image]
\addplot[domain=0:40.8,samples=60,very thick,accent]
  ({x},{tan(45)*x-9.8*x^2/(2*400*cos(45)^2)});
% velocity arrows: launch, apex, descent (components)
\draw[->,very thick] (axis cs:0,0) -- (axis cs:7.07,7.07) node[midway,left,font=\footnotesize] {$\mathbf v_0$};
\draw[->,thick,dashed] (axis cs:0,0) -- (axis cs:7.07,0) node[pos=0.85,below,font=\footnotesize] {$v_x$};
\draw[->,thick,dashed] (axis cs:7.07,0) -- (axis cs:7.07,7.07) node[midway,right,font=\footnotesize] {$v_{y0}$};
\draw[->,very thick] (axis cs:20.4,10.2) -- (axis cs:27.5,10.2) node[midway,above,font=\footnotesize] {$\mathbf v$: pure $v_x$};
\draw[->,very thick] (axis cs:32,6.3) -- (axis cs:39,1.4) node[midway,right,font=\footnotesize] {$\mathbf v$};
\draw[->,thick,dashed] (axis cs:32,6.3) -- (axis cs:39,6.3);
\draw[->,thick,dashed] (axis cs:39,6.3) -- (axis cs:39,1.4) node[midway,right,font=\footnotesize] {$v_y<0$};
\fill (axis cs:20.4,10.2) circle (2pt) node[above=2pt,font=\footnotesize] {apex};
\end{axis}
\end{tikzpicture}

\caption{Projectile with $v_0=\qty{20}{m/s}$ at $45^{\circ}$. Notice what
changes along the path: $v_x$ never does; $v_y$ drains to zero at the apex and
goes negative after. The parabola is a parametric curve --- $y(x)$ has no
time in it; recover $t$ from $x=v_{0x}t$ whenever it is needed.}
\label{fig:proj}
\end{figure}

The flat-ground range $R=\frac{v_0^{2}\sin2\theta}{g}$ follows by setting
$y=0$ and eliminating $t$. The generalised version --- landing at height $h$
--- is worth more:

\begin{example}[Landing on a ledge]\label{ex:elevated}
The launch of figure~\ref{fig:proj} ($v_0=\qty{20}{m/s}$, $45^{\circ}$)
lands on a platform $h=\qty{5}{m}$ above the launch level. Find the range.
\par\smallskip
\textbf{Strategy.} Do not memorise a second range formula: solve
$y(t)=h$ for the \emph{later} crossing time, then $R=v_{0x}t$.
\begin{solution}
$v_{0x}=v_{0y}=20/\sqrt2=14.14\,\mathrm{m/s}$. Solve
$14.14t-4.9t^{2}=5$, i.e.\ $4.9t^{2}-14.14t+5=0$:
$t=\frac{14.14\pm\sqrt{199.9-98}}{9.8}=\frac{14.14\pm10.10}{9.8}$; the
later root $t\approx2.47\,\mathrm{s}$ (the earlier one,
$0.41\,\mathrm{s}$, is the upward crossing). Then
$R=14.14\times2.47\approx\qty{34.9}{m}$ --- \emph{shorter} than the
flat-ground $40.8\,\mathrm{m}$, because the ledge intercepts the
descending arc early. Landing $5\,\mathrm{m}$ \emph{below} gives
$R\approx45.3\,\mathrm{m}$; the simulator prints both signs of $h$ in
section~\ref{sec:kin-code}. \checkmark
\end{solution}
\end{example}

\section{Circular motion: acceleration without speeding up}\label{sec:kin-circ}

Uniform circular motion --- constant speed $v$ on a circle of radius $r$ ---
is accelerated: the velocity vector turns even though its length never
changes. Figure~\ref{fig:centripetal} draws the velocity at two nearby
instants: $\mathbf v(t)$ and $\mathbf v(t+\Delta t)$ are equal-length vectors
differing by a small rotation, so $\Delta\mathbf v$ points \emph{inward}, and
in the limit $\Delta t\to0$ it points exactly at the centre. The magnitude
comes from arc geometry: the velocity vector turns through the same angle the
radius does, $\Delta\theta=v\Delta t/r$, so
$|\Delta\mathbf v|\approx v\,\Delta\theta=v^{2}\Delta t/r$, giving
\begin{equation}
  a_c=\frac{v^{2}}{r}\quad\text{toward the centre.}\label{eq:ac}
\end{equation}
This is \emph{centripetal acceleration}\index{centripetal acceleration} ---
required by geometry, supplied by whatever real force (tension, friction,
gravity) is handy; chapter~2 matches it to its source.

\begin{figure}[htbp]\centering
% dv geometry for centripetal acceleration
\begin{tikzpicture}[scale=1.1,>=stealth]
  % left: circle with two velocity vectors
  \begin{scope}
    \draw[thick,gray] (0,0) circle (1.9);
    \draw[->,very thick,accent] (1.9,0) -- ++(0,1.7) node[right,font=\small] {$\mathbf v(t)$};
    \draw[->,very thick,accent] ({1.9*cos(40)},{1.9*sin(40)}) -- ++({-1.7*sin(40)},{1.7*cos(40)}) node[above left=-2pt,font=\small] {$\mathbf v(t+\Delta t)$};
    \fill (0,0) circle (1.5pt) node[below left,font=\small] {$O$};
    \draw[dashed] (1.9,0) -- ({1.9*cos(40)},{1.9*sin(40)});
    \node[font=\footnotesize] at (0.7,-2.4) {speed $|\mathbf v|$ constant};
  \end{scope}
  % right: the velocity difference triangle
  \begin{scope}[xshift=6.4cm,yshift=-0.6cm]
    \draw[->,very thick,accent] (0,0) -- (0,1.9) node[midway,left,font=\small] {$\mathbf v(t)$};
    \draw[->,very thick] (0,0) -- (-0.85,1.85) node[midway,left,font=\small] {$\mathbf v(t+\Delta t)$};
    \draw[->,very thick] (0,1.9) -- (-0.85,1.85) node[midway,above,font=\small] {$\Delta\mathbf v$};
    \node[font=\footnotesize,align=center] at (-0.2,-0.5) {$\Delta\mathbf v$ points inward:\\as $\Delta t\to0$, toward $O$};
  \end{scope}
\end{tikzpicture}

\caption{Where $a_c=v^{2}/r$ comes from. Notice the argument needs only that
$|\mathbf v|$ is constant: the velocity arrow rotates with the radius, so
$\Delta\mathbf v$ is perpendicular to $\mathbf v$ and aimed at the centre.}
\label{fig:centripetal}
\end{figure}

\begin{example}[A car on a curve]\label{ex:circular}
A car takes a radius-$50\,\mathrm{m}$ curve at $20\,\mathrm{m/s}$
($72\,\mathrm{km/h}$). How large is the sideways acceleration, in units of
$g$?
\par\smallskip
\textbf{Strategy.} Equation \eqref{eq:ac} directly; compare with $g$ for
intuition (your body feels acceleration in $g$'s).
\begin{solution}
$a_c=20^{2}/50=8\,\mathrm{m/s^{2}}\approx0.8g$. The tires must supply every
newton of it --- on ice, $20\,\mathrm{m/s}$ around this curve is impossible,
which is why ramps are banked (chapter~2). \checkmark
\end{solution}
\end{example}

\begin{exercise}[medium]
A CD spins at $500\,\mathrm{rpm}$. Find $a_c$ at its $6\,\mathrm{cm}$ rim in
SI and in units of $g$. (AP-C)
\end{exercise}

\section{Relative velocity and Galilean frames}\label{sec:kin-rel}

All velocities are relative to a frame; switching frames adds a constant
vector:
\[
  \mathbf v_{A/C}=\mathbf v_{A/B}+\mathbf v_{B/C},
\]
the \emph{Galilean velocity addition}\index{relative velocity} --- read it as
``$A$ as seen from $C$ equals $A$ as seen from $B$ plus $B$ as seen from
$C$''. The subscripts chain like a telescope. Figure~\ref{fig:reltri} is the
one picture to carry: two perpendicular legs and a resultant.

\begin{figure}[htbp]\centering
% relative velocity triangle: boat + river = ground velocity
\begin{tikzpicture}[scale=0.9,>=stealth]
  % river current
  \draw[->,very thick] (0,0) -- (3,0) node[midway,below,font=\small] {$\mathbf v_{\mathrm{river}}=\qty{3}{m/s}$};
  % boat across (its own thrust)
  \draw[->,very thick,accent] (0,0) -- (0,2.4) node[midway,left,font=\small] {$\mathbf v_{\mathrm{boat/river}}=\qty{4}{m/s}$};
  % resultant
  \draw[->,very thick,dashed] (0,0) -- (3,2.4) node[midway,above left,font=\small] {$\mathbf v_{\mathrm{boat/ground}}=\qty{5}{m/s}$};
  \draw[dotted,gray] (0,2.4) -- (3,2.4);
  \draw[dotted,gray] (3,0) -- (3,2.4);
  % right-angle mark
  \draw (0.4,0) -- (0.4,0.4) -- (0,0.4);
  \node[font=\footnotesize,align=center] at (4.4,1.2) {$v_{\mathrm{bg}}=v_{\mathrm{br}}+v_{\mathrm{rg}}$\\ Galilean addition};
\end{tikzpicture}

\caption{Relative velocity as vector addition. Notice the frame
$B=\mathrm{river}$ does all the work: the boat's thrust is measured relative
to the water it pushes against, and the bank sees the sum.}
\label{fig:reltri}
\end{figure}

\begin{example}[River crossing]\label{ex:river}
A boat that makes $4\,\mathrm{m/s}$ in still water aims straight across a
$200\,\mathrm{m}$ river flowing at $3\,\mathrm{m/s}$. When does it land, and
where? Then: what heading would land it straight across?
\par\smallskip
\textbf{Strategy.} Components decouple, exactly like projectile motion:
crossing time depends on the across-component only. For the straight
crossing, aim upstream so the upstream component of boat velocity cancels
the current.
\begin{solution}
Across-component $4\,\mathrm{m/s}$: crossing time $t=200/4=50\,\mathrm{s}$.
Drift $=3\times50=150\,\mathrm{m}$ downstream --- a
$3{-}4{-}5$ triangle, ground speed $5\,\mathrm{m/s}$. To land straight
across, aim at angle $\theta$ upstream with $4\sin\theta=3$, i.e.\
$\theta=\arcsin0.75\approx48.6^{\circ}$; the effective across-speed drops
to $4\cos\theta=\sqrt{7}\approx2.65\,\mathrm{m/s}$, so
$t\approx75.7\,\mathrm{s}$ --- straighter costs slower. \checkmark
\end{solution}
\end{example}

\section{The simulator: integration is not optional}\label{sec:kin-code}

Every claim above is checkable by \emph{numerically} running the calculus
loop --- and in real engineering (drag, control forces, planetary $n$-body)
numerical integration is the only option. The listing below integrates
projectile motion by forward Euler, measures the step-size error, then adds
quadratic drag --- the first preview of chapter 2's terminal velocity.

\begin{checkpoint}[Predict before you run]
Euler overshoots: $v$ is updated after $x$, so each step uses a stale
velocity. Guess the sign of the range error before looking.
\end{checkpoint}

\begin{lstlisting}[language=Python,caption={\texttt{code/ch01\_projectile.py} --- Euler integration vs closed form, then with quadratic drag.},label={lst:proj}]
def euler_range(v0, th, dt, k=0.0):
    x = y = 0.0
    vx, vy = v0*math.cos(th), v0*math.sin(th)
    while y >= -1e-12:
        s = math.hypot(vx, vy)
        ax, ay = -k*s*vx, -G - k*s*vy   # gravity + quadratic drag
        x += vx*dt;  y += vy*dt       # update position first
        vx += ax*dt; vy += ay*dt      # then velocity: forward Euler
    return x
\end{lstlisting}

Output of a real run:
\begin{lstlisting}[numbers=none]
analytic range:  40.8163 m
euler dt=0.1    range= 42.4264 m  error=+1.6101 m
euler dt=0.01   range= 41.0122 m  error=+0.1959 m
euler dt=0.001  range= 40.8425 m  error=+0.0262 m
with quadratic drag k=0.02: range=25.8593 m (36.6% shorter)
landing at h= +5.0 m: range=34.9825 m (vs 40.8163 on flat ground)
landing at h= -5.0 m: range=45.3195 m (vs 40.8163 on flat ground)
\end{lstlisting}

Notice two things. First, the error shrinks roughly with $h$ --- Euler is
first-order --- and overshoots positive, exactly as predicted: the stale
velocity keeps climbing after it should turn. Second, drag at
$k=0.02\,\mathrm{m^{-1}}$ cuts the range by a third: real ballistics is not
the parabola, and chapter 2's drag force is why.

\section{In engineering practice: dead reckoning}\label{sec:kin-eng}

An inertial measurement unit outputs acceleration, never position --- so
every drone, submarine, and phone navigation chip lives inside the right
half of figure~\ref{fig:loop}: $v(t)=v_0+\int a$, $x(t)=x_0+\int v$, run
thousands of times a second on a tiny processor. The cost is
\emph{drift}\index{drift (IMU)}: each integration accumulates the sensor's
bias error, and a $\qty{1}{mg}$ accelerometer bias becomes a
$\sim 0.5\,\mathrm{m}$ position error in ten seconds. GPS exists to reset
that integral: satellite fixes pin $x(t)$ periodically so the loop restarts
from truth. Dead reckoning --- this chapter, continuously --- is the
engineering version of ``kinematics is the calculus loop.''

\section{Common mistakes}\label{sec:kin-mistakes}

\begin{tripwire}[Speed $\neq$ velocity]
Distance travelled and displacement differ the moment the path bends or
reverses. ``How fast'' questions can mean speed or velocity --- check the
sign convention of the problem before answering.
\end{tripwire}
\begin{tripwire}[$a_c=v^{2}/r$ is not a new force]
It is a \emph{required} acceleration. The forces are tension/friction/
gravity; their sum must equal $mv^{2}/r$ inward, not the other way around.
\end{tripwire}
\begin{tripwire}[Inconsistent sign of $g$]
Pick an axis once. With $+y$ up, $a=-9.8$; with $+y$ down, $a=+9.8$. Mixing
both inside one solution is the most common source of wrong answers in
$\S\ref{sec:kin-consta}$--$\ref{sec:kin-proj}$.
\end{tripwire}
\begin{tripwire}[Constant-$a$ formulas on varying $a$]
Equations \eqref{eq:consta1}--\eqref{eq:consta2} need $a=\text{const}$. With
drag, springs, or powered flight they are wrong; integrate the loop instead
(section~\ref{sec:kin-code}).
\end{tripwire}
\begin{tripwire}[Adding velocities as scalars]
$\mathbf v_{A/C}=\mathbf v_{A/B}+\mathbf v_{B/C}$ is vector addition. A
$4$-m/s boat in a $3$-m/s river does \emph{not} move at $7\,\mathrm{m/s}$ ---
it moves at $5\,\mathrm{m/s}$ diagonally (figure~\ref{fig:reltri}).
\end{tripwire}

\section{Summary}\label{sec:kin-summary}

\begin{table}[ht]\centering
\begin{tabularx}{\linewidth}{@{}lY@{}}
\toprule
\tablehead\thead{Idea} & \thead{One-line version}\\
\midrule
calculus loop & $x\xrightarrow{d/dt}v\xrightarrow{d/dt}a$; integrate back, paying two constants \\
constant $a$ & $v=v_0+at$, $x=x_0+v_0t+\tfrac12at^{2}$, $v^{2}=v_0^{2}+2a\Delta x$ \\
projectile & $x$: const.\ velocity; $y$: const.\ accel $-g$; shared only by $t$ \\
range on height $h$ & solve $y(t)=h$, take the later root, $R=v_{0x}t$ \\
circular motion & $a_c=v^{2}/r$ toward centre --- required, not a force \\
relative velocity & $\mathbf v_{A/C}=\mathbf v_{A/B}+\mathbf v_{B/C}$ (vectors) \\
speed rule & $v,a$ same sign $\Rightarrow$ speeding up; opposite $\Rightarrow$ slowing (fig.~\ref{fig:signmap}) \\
\bottomrule
\end{tabularx}
\caption{Kinematics cheat sheet.}
\label{tab:kinsummary}
\end{table}

\begin{figure}[htbp]\centering
% sign map: speed change from (v,a) signs
\begin{tikzpicture}[scale=1.0]
  \draw[thick] (-0.1,-0.1) rectangle (4.6,3.4);
  \draw[thick] (2.25,-0.1) -- (2.25,3.4);
  \draw[thick] (-0.1,1.65) -- (4.6,1.65);
  \node[font=\footnotesize] at (4.0,2.9) {$a>0$};
  \node[font=\footnotesize] at (4.0,0.4) {$a<0$};
  \node[font=\footnotesize,rotate=90] at (0.2,2.65) {$v>0$};
  \node[font=\footnotesize,rotate=90] at (0.2,0.65) {$v<0$};
  \node[font=\small,align=center,fill=accentsoft] at (1.1,2.5) {speeding up};
  \node[font=\small,align=center] at (3.4,2.5) {slowing down};
  \node[font=\small,align=center] at (1.1,0.75) {slowing down};
  \node[font=\small,align=center,fill=accentsoft] at (3.4,0.75) {speeding up};
  \node[font=\footnotesize,align=center] at (2.25,-0.6) {$v$ and $a$ \emph{same} sign $\Rightarrow$ faster; opposite $\Rightarrow$ slower};
\end{tikzpicture}

\caption{The $(v,a)$ sign map. Notice it never asks about position: only the
\emph{relative} signs of velocity and acceleration set whether the object
speeds up or slows down.}
\label{fig:signmap}
\end{figure}

\begin{wins}[Definition of done]
You can now: read $v$ and $a$ off an $x(t)$ graph and off a data table;
derive the three constant-acceleration formulas; split projectile motion into
components and land it on a ledge; explain where $a_c=v^{2}/r$ comes from;
add velocities through a frame chain; and run the calculus loop numerically
with a first-order integrator.
\end{wins}

\chapternotes{%
The calculus-first ordering and the derivations follow MIT 8.01SC weeks 1--3
\cite{ocw801}; AP Physics C unit 1 covers the same ground in the MC/FR split
\cite{apced}. OpenStax College Physics chapters 3--4 \cite{openstax1} supply
the projectile and relative-velocity treatment; the centripetal derivation by
velocity difference is Feynman I \S8-5 \cite{feynman1}. Halliday--Resnick
\cite{hrw} has the largest drill inventory; Morin \cite{morin} is the harder
recommendation for anyone who found this chapter easy. The dead-reckoning
paragraph sketches the strapdown inertial navigation problem; the bias-drift
scaling is standard.}

\begin{rulebox}[{Watch instead}]
\begin{enumerate}
\item \emph{1D Kinematics --- Speed, Velocity, Acceleration} --- Walter Lewin,
MIT 8.01x Lecture 2, 50:42, \url{https://youtu.be/q9IWoQ199_o} ---
\S\ref{sec:kin-pos}--\S\ref{sec:kin-consta} at full length, same notation
(accessed 2026-10-11).
\item \emph{Why distance is area under the velocity--time line} --- Khan
Academy, 9:25, \url{https://youtu.be/d-_eqgj5-K8} --- the integration half of
\S\ref{sec:kin-acc} in eight minutes (accessed 2026-10-11).
\end{enumerate}
\end{rulebox}

\section*{Further reading}
\addcontentsline{toc}{section}{Further reading}

\begin{itemize}
\item MIT OpenCourseWare, \emph{8.01SC Classical Mechanics}, weeks 1--3 ---
the source course; recitation videos solve the same problems. Free at
\url{https://ocw.mit.edu}, course 8.01SC'' (verified 2026-10-11).
\item OpenStax, \emph{College Physics 2e}, ch.~3 ``Two-Dimensional
Kinematics'' --- free text, more worked projectile examples. At
\url{https://openstax.org}, search ``College Physics 2e'' (verified 2026-10-11).
\item College Board, \emph{AP Physics C: Mechanics CED}, Unit 1 --- the exam
blueprint this chapter aligns to; FR verb practice. At
\url{https://apstudents.collegeboard.org}, course page ``AP Physics C:
Mechanics'' (verified 2026-10-11).
\item R.~P. Feynman, \emph{The Feynman Lectures on Physics}, vol.~I, ch.~8
``Motion'' --- the prose version, still the best motivation for why
derivatives run physics. Free online at Caltech.
\item D.~Morin, \emph{Introduction to Classical Mechanics}, ch.~2 ---
kinematics for people who want the hard version; ISBN 978-0521876223.
\item Khan Academy, \emph{Physics --- one-dimensional motion} --- drill
exercises matching \S\ref{sec:kin-pos}--\S\ref{sec:kin-consta}, at
\url{https://www.khanacademy.org/science/physics}. Verified 2026-10-11.
\item The Physics Classroom, \emph{1-D and 2-D Kinematics} tutorials ---
readable remediation with check-yourself quizzes, at
\url{https://www.physicsclassroom.com} (verified 2026-10-11).
\end{itemize}

\section*{Additional exercises}
\addcontentsline{toc}{section}{Additional exercises}

\begin{exercise}[easy]
From the $v(t)$ graph of figure~\ref{fig:triple}: read off the launch
velocity, the time of the peak, and the landing time --- then confirm all
three from $x(t)=20t-4.9t^{2}$. (AP-C)
\end{exercise}
\begin{exercise}[easy]
Convert these accelerations to SI: $32\,\mathrm{km/h/s}$,
$0.05\,\mathrm{miles/h^{2}}$ (in $\mathrm{m/s^{2}}$, approximately), and
$9.8\,\mathrm{m/s^{2}}$ in $\mathrm{km/h^{2}}$. (AP-C)
\end{exercise}
\begin{exercise}[easy]
A ball is dropped from rest at height $h$. Find its speed after $t$ seconds
and after falling $d$ metres --- which of the three constant-$a$ formulas did
each answer need? (F=ma)
\end{exercise}
\begin{exercise}[easy]
A train accelerates uniformly from rest to $30\,\mathrm{m/s}$ over
$600\,\mathrm{m}$. Find $a$ and the elapsed time. (AP-C)
\end{exercise}
\begin{exercise}[easy]
Displacement vs.\ distance: a runner does one lap of a $400\,\mathrm{m}$
track in $80\,\mathrm{s}$. Compute average velocity and average speed. (AP-C)
\end{exercise}
\begin{exercise}[medium]
Show that for a launch on level ground the maximum range at fixed $v_0$
comes at $\theta=45^{\circ}$, then show the max range equals
$\frac{v_0^{2}}{g}$ --- by maximising $R(\theta)$, not by quoting it. (AP-C FR-style)
\end{exercise}
\begin{exercise}[medium]
A second crossing: the boat of example~\ref{ex:river} wants minimum
\emph{time} rather than zero drift. Show the fastest crossing points
straight across ($t=50\,\mathrm{s}$), and compute where it lands. Then find
the heading for minimum \emph{total distance travelled over ground}.
(F=ma)
\end{exercise}
\begin{exercise}[medium]
From the dragster table of example~\ref{ex:dataderiv}, estimate $v(2.0)$
and $a(2.0)$ by the same symmetric-difference method, then fit
$x(t)\approx\tfrac12at^{2}$ and compare. (F=ma, data analysis)
\end{exercise}
\begin{exercise}[medium]
A ball is thrown horizontally at $15\,\mathrm{m/s}$ off a
$30\,\mathrm{m}$ cliff. Find the time to land, the range, and the velocity
vector (magnitude and direction) at impact. (AP-C)
\end{exercise}
\begin{exercise}[medium]
On the sign map (figure~\ref{fig:signmap}): an object has $v<0$ and $a<0$.
Is it speeding up? Give the everyday example of each quadrant.
(AP-C)
\end{exercise}
\begin{exercise}[medium]
An airplane flies at $200\,\mathrm{m/s}$ north relative to air, in a
$40\,\mathrm{m/s}$ east wind. Find the ground track (heading and ground
speed). (F=ma)
\end{exercise}
\begin{exercise}[medium]
Two cars start $500\,\mathrm{m}$ apart and drive toward each other at
$25\,\mathrm{m/s}$ and $20\,\mathrm{m/s}$. When and where do they meet ---
solve in the ground frame, then in the frame of the slower car. (F=ma)
\end{exercise}
\begin{exercise}[hard]
Generalise example~\ref{ex:elevated}: for a launch at $(v_0,\theta)$ onto a
ledge at height $h$, find the angle $\theta^{*}$ that maximises range. Show
$\theta^{*}\to45^{\circ}$ as $h\to0$ and find the leading correction for
small $h/v_0^{2}$. (USAPhO-style)
\end{exercise}
\begin{exercise}[hard]
A pursuit with a cap: a cheetah accelerates at $5\,\mathrm{m/s^{2}}$ but is
speed-limited to $30\,\mathrm{m/s}$; a gazelle, $20\,\mathrm{m}$ ahead,
accelerates at $3\,\mathrm{m/s^{2}}$ capped at $25\,\mathrm{m/s}$. When does
the cheetah catch it? (Track the phase where both are at cap.) (APhO-style)
\end{exercise}
\begin{exercise}[hard]
From first principles: derive $a_c=v^{2}/r$ without the velocity-difference
argument --- write $\mathbf r(t)=(r\cos\omega t,\,r\sin\omega t)$ and
differentiate twice. Confirm both methods agree on direction and magnitude.
(USAPhO-style)
\end{exercise}
